\documentclass[12pt, a4paper]{article}
\usepackage[utf8]{inputenc}
\usepackage[T2A]{fontenc}
\usepackage{indentfirst, setspace}
\usepackage{tabularx, multirow}
\usepackage[normalem]{ulem}
\usepackage[style=russian]{csquotes}
\usepackage[english,russian]{babel}
\usepackage{hyperref}
\usepackage{ragged2e}
\usepackage{caption}
\usepackage{wrapfig}
\usepackage{amsmath}
\usepackage{tikz}
\makeatletter
\def\@biblabel#1{#1. }
\makeatother
\captionsetup{labelsep=endash}
\usepackage{listings}
\linespread{1.3}
\lstset{
  language=C++, % Язык программирования (можно выбрать другой)
  basicstyle=\ttfamily\small, % Основной стиль шрифта
  keywordstyle=\color{blue}, % Стиль ключевых слов
  breaklines=true,
  commentstyle=\color{green}, % Стиль комментариев
  stringstyle=\color{red}, % Стиль строк
  extendedchars=\true,
  showstringspaces=false,
}
\usepackage[left=20mm, top=15mm, right=15mm, bottom=20mm]{geometry}

\begin{document}
\begin{titlepage}
     \fontsize{12}{12}\selectfont

  {\centering

   \begin{bf}

    \begin{wrapfigure}{l}{19mm}
        
\includegraphics[width=22mm]{emblem.jpeg}
    \end{wrapfigure}


    \noindent Министерство науки и высшего образования Российской Федерации

    \noindent Федеральное государственное бюджетное образовательное учреждение высшего образования

    \noindent \enquote{Московский государственный технический унверситет

     \noindent имени Н.Э. Баумана

     \noindent (национальный исследовательский университет)}

    \noindent (МГТУ им. Н.Э. Баумана)

   \end{bf}
  }

  \vspace{0.4cm}

  {\setstretch{0.1}
   \noindent\rule{\textwidth}{1mm}
   \noindent\rule{\textwidth}{0.5mm}

  }

  \fontsize{14}{21}\selectfont

  \noindent\begin{tabularx}{\textwidth}{l >{\centering\arraybackslash}X}
   ФАКУЛЬТЕТ & \flqq Фундаментальные Науки\frqq \\ \cline{2-2}

   КАФЕДРА & ФН-12 \flqq Математическое моделирование\frqq \\ \cline{2-2}
  \end{tabularx}


  \vspace{1cm}


  \begin{center}
   \begin{bf}

    \fontsize{24}{36}\selectfont
    ОТЧЕТ

    \fontsize{20}{30}\selectfont
    ПО ЛАБОРАТОРНОЙ РАБОТЕ НА ТЕМУ:

    Перемножение матриц

   \end{bf}
  \end{center}

  \fontsize{14}{21}\selectfont
  \vspace{5cm}


  \noindent\begin{tabularx}{\textwidth}{ X >{\centering}p{4cm} p{1cm} c }
   Студент: & & & Подобедов В. В. \\ \cline{2-2} \cline{4-4}
   & \fontsize{10}{15}\selectfont дата, подпись & & \fontsize{10}{15}\selectfont Ф.И.О. \\
   Преподаватель: & & & Волкова Л. Л.\\ \cline{2-2} \cline{4-4}
   & \fontsize{10}{15}\selectfont дата, подпись & & \fontsize{10}{15}\selectfont Ф.И.О.
  \end{tabularx}

  \vspace{\fill}

  \begin{center}
   \it{Москва}, 2023
  \end{center}

  \thispagestyle{empty}
\end{titlepage}

\renewcommand{\contentsname}{Содержание}
\setcounter{page}{2}
\tableofcontents

\newpage
\begin{center}
   \section*{Введение} 
\addcontentsline{toc}{section}{Введение}
\end{center}
\justifying
\textbf{Цель лабораторной работы}: провести сравнительный анализ алгоритмов умножения матриц.

Для достижения поставленной цели требуется решить следующие \textbf{задачи}.
\begin{enumerate}
\item Описать три алгоритма умножения матриц --- классический, Винограда и Штрассена.
\item Реализовать разработанные алгоритмы.
\item Выполнить тестирование реализации разработанных алгоритмов.
\item Дать оценку сложности алгоритмов.
\item Оценить эффективность каждого алгоритма (по памяти и времени).
\end{enumerate}

\newpage
\begin{center}
   \section{Аналитическая часть} 
\end{center}

\textbf{Матрица} --- математический объект, записываемый в виде прямоугольной таблицы, который представляет собой совокупность строк и столбцов, на пересечении которых находятся его элементы. Количество строк и столбцов задает размер матрицы.

\textbf{Классический алгоритм умножения матриц} --- это стандартный метод перемножения матриц в линейной алгебре. Если есть две матрицы $A$ размером $m \times n$ и $B$ размером $n \times p$, то произведение матриц $C = AB$ будет матрицей размером $m \times p$. Каждый элемент $C[i][j]$ матрицы C вычисляется путем умножения элементов соответствующей строки матрицы $A$ на соответствующие элементы столбца матрицы $B$ и последующим сложением результатов. 

\textbf{Алгоритм Винограда для умножения матриц}~\cite{book} --- это метод, где предварительно вычисляются некоторые значения и используются для уменьшения количества умножений в основном цикле умножения матриц. Этот метод полезен для больших матриц, где вычислительные затраты могут быть значительными. 

\textbf{Алгоритм Штрассена для умножения матриц}~\cite{book} --- это метод умножения матриц, который уменьшает количество операций умножения стандартным способом. Он основан на разделении матриц на подматрицы меньшего размера, что позволяет снизить сложность умножения. Основная идея алгоритма заключается в том, что вместо классического умножения матриц $A$ и $B$ размерности $n \times n$ можно разбить каждую из них на четыре подматрицы размера $\frac{n}{2} \times \frac{n}{2}$, что дает $8$ уравнений для вычисления произведения. Однако алгоритм Штрассена использует не все эти уравнения, а только $7$, чтобы сократить количество операций умножения.

\newpage
\begin{center}
    \section{Конструкторская часть}
\end{center}
\subsection{Умножение матриц классическим способом}
Формула для вычисления элемента $C[i][j]$ выглядит следующим образом~(\ref{eq:classic}):
\begin{equation} \label{eq:classic}
C_{ij} = \sum_{k=1}^{n} A_{ik} \cdot B_{kj}
\end{equation}

Нетрудно понять, что в данной формуле $i$ --- номер строки, $j$ --- номер столбца, $k$ --- индекс суммирования, $C$ --- резульирующая матрица, $A$ --- первая матрица, $B$ --- вторая матрица. Это означает, что каждый элемент $C[i, j]$ матрицы C равен сумме произведений соответствующих элементов строки $i$ матрицы $A$ и столбца $j$ матрицы $B$. Процесс вычисления всех элементов матрицы $C$ повторяется для каждой пары $(i, j)$ в пределах размерности матрицы C.

\subsection{Метод Винограда}
Алгоритм Винограда включает следующие шаги.
\begin{enumerate}
    \item Создать новую матрицу $C$ размерностью $m \times n$ и заполните ее нулями.
    \item Создать вспомогательные массивы $row\_factor$ и $col\_factor$ для оптимизации вычислений. Их размерность равна числу строк матрицы $A$ и числу столбцов матрицы $B$ соответственно.
    \item Вычислить $row\_factor[i]$ как сумму произведений пар элементов из $A[i]$, умноженных на соответствующие элементы из $A[i]$.
    \item Вычислите $col\_factor[j]$ как сумму произведений пар элементов из $B[j]$, умноженных на соответствующие элементы из $B[j]$.
    \item Используя предварительно вычисленные $row\_factor$ и $col\_factor$, вычислите каждый элемент матрицы $C$ следующим образом: 
    \[ C[i][j] = rc\_f + \sum_{k=0}^{k-1} (A[i][k] + B[k][j]) \cdot (A[i][k+1] + B[k+1][j]) \text{, где}\]  \[ rc\_f = -row\_factor[i] - col\_factor[j] \] 
    \item Если размерность матрицы нечетна, нужно добавить к соответствующему элементу $C$ произведение соответствующих элементов из $A$ и $B$, которые не участвовали в основных вычислениях.
\end{enumerate}

\subsection{Метод Штрассена}
Шаги алгоритма Штрассена, следующие.
\begin{enumerate}
    \item Каждая изматриц $A$ и $B$ разбивается на четыре части размера $\frac{n}{2} \times \frac{n}{2}$. Таким образом получаются такие матрицы: 
    $\begin{pmatrix}
      A_{11} & A_{12}\\
      A_{21} & A_{22}
    \end{pmatrix}$
    и 
    $\begin{pmatrix}
      B_{11} & B_{12}\\
      B_{21} & B_{22}
    \end{pmatrix}$
    \item Вычисляются промежуточные матрицы от $M_1$ до $M_7$ по формулам:
    \begin{align*}
      M_1 &= (A_{11} + A_{22}) \cdot (B_{11} + B_{22}) \\
      M_2 &= (A_{21} + A_{22}) \cdot B_{11} \\
      M_3 &= A_{11} \cdot (B_{12} - B_{22}) \\
      M_4 &= A_{22} \cdot (B_{21} - B_{11}) \\
      M_5 &= (A_{11} + A_{12}) \cdot B_{22} \\
      M_6 &= (A_{21} - A_{11}) \cdot (B_{11} + B_{12}) \\ 
      M_7 &= (A_{12} - A_{22}) \cdot (B_{21} + B_{22})
    \end{align*}
    \item Используя промежуточные матрицы $M_1$ до $M_7$, вычисляются четыре части результирующей матрицы $C$:
    \begin{align*}
      C_{11} &= M_1 + M_4 - M_5 + M_7 \\
      C_{12} &= M_3 + M_5 \\
      C_{21} &= M_2 + M_4 \\
      C_{22} &= M_1 - M_2 + M_3 + M_6
    \end{align*}
    \item Полученные части результирующей матрицы $C_{11}$, $C_{12}$, $C_{21}$, $C_{22}$ объединяются в единую матрицу $C$.
\end{enumerate}

\subsection{Функции программы}
Требуется разработать следующие функции.
\begin{itemize}
    \item[---] Функция ``matrixAddition'' --- функция, которая складывает марицы.
    \item[---] Функция ``matrixSubtraction'' --- функция, которая перемножает две матрицы.
    \item[---] Функция ``strassenMatrixMultiply'' --- функция для умножения матриц методом Штрассена.
    \item[---] Функция ``classicMatrixMultiply'' --- функция для умножения матриц классическим способом.
    \item[---] Функция ``matrixMultiplyVinograd'' ---  Функция для умножения матриц с использованием алгоритма Винограда.
    \item[---] Функция ``printMatrix'' --- функция для вывода матрицы.
\end{itemize}

К программе предъявлены следующие требования. В начале выполнения работы программы пользователь выбирает каким методом будут перемножены матрицы. Далее с помощью выбранного алгоритма вычисляется произведение матриц и на экран консоли выводится результирующая матрица.

\newpage
\begin{center}
    \section{Технологическая часть}
\end{center}

Лабораторная работа была реализована на языке программирования C++. Кроме стандартной библиотеки <iostream>, которая предоставляет возможность пользоваься средствами ввода-вывода для стандартной консоли, была подключена библиотека <vector>. Она включает в себя шаблон класса для контейнеров последовательности. Вектор хранит элементы заданного типа в линейном расположении и обеспечивает быстрый случайный доступ к любому элементу. Также для вычисления работы программы была взаимствована библиотека <ctime>.

Далее приведен листинг реализации трех алгоритмов (листинг~\ref{lst:label1}).
\begin{lstlisting}[caption={Исходный код программы для нахождения произведения матриц разными методами},label={lst:label1}]
#include <iostream>
#include <vector>
#include <ctime>

using namespace std;

// Функция для сложения матриц
vector<vector<int>> matrixAddition(const vector<vector<int>>& A, const vector<vector<int>>& B) {
	int n = A.size();
	vector<vector<int>> result(n, vector<int>(n, 0));

	for (int i = 0; i < n; i++) {
		for (int j = 0; j < n; j++) {
			result[i][j] = A[i][j] + B[i][j];
		}
	}

	return result;
}

// Функция для вычитания матриц
vector<vector<int>> matrixSubtraction(const vector<vector<int>>& A, const vector<vector<int>>& B) {
	int n = A.size();
	vector<vector<int>> result(n, vector<int>(n, 0));
	for (int i = 0; i < n; i++) {
		for (int j = 0; j < n; j++) {
			result[i][j] = A[i][j] - B[i][j];
		}
	}

	return result;
}

// Функция для умножения матриц методом Штрассена
vector<vector<int>> strassenMatrixMultiply(const vector<vector<int>>& A, const vector<vector<int>>& B) {
	int n = A.size();

	if (n <= 2) {
		vector<vector<int>> result(n, vector<int>(n, 0));
		for (int i = 0; i < n; i++) {
			for (int j = 0; j < n; j++) {
				for (int k = 0; k < n; k++) {
					result[i][j] += A[i][k] * B[k][j];
				}
			}
		}

		return result;
	}

	int newSize = n;
	if (n % 2 != 0) {
		newSize++;
	}

	vector<vector<int>> newA(newSize, vector<int>(newSize, 0));
	vector<vector<int>> newB(newSize, vector<int>(newSize, 0));
	for (int i = 0; i < n; i++) {
		for (int j = 0; j < n; j++) {
			newA[i][j] = A[i][j];
			newB[i][j] = B[i][j];
		}
	}

	vector<vector<int>> A11(newSize / 2, vector<int>(newSize / 2, 0));
	vector<vector<int>> A12(newSize / 2, vector<int>(newSize / 2, 0));
	vector<vector<int>> A21(newSize / 2, vector<int>(newSize / 2, 0));
	vector<vector<int>> A22(newSize / 2, vector<int>(newSize / 2, 0));
	vector<vector<int>> B11(newSize / 2, vector<int>(newSize / 2, 0));
	vector<vector<int>> B12(newSize / 2, vector<int>(newSize / 2, 0));
	vector<vector<int>> B21(newSize / 2, vector<int>(newSize / 2, 0));
	vector<vector<int>> B22(newSize / 2, vector<int>(newSize / 2, 0));
	for (int i = 0; i < newSize / 2; i++) {
		for (int j = 0; j < newSize / 2; j++) {
			A11[i][j] = newA[i][j];
			A12[i][j] = newA[i][j + newSize / 2];
			A21[i][j] = newA[i + newSize / 2][j];
			A22[i][j] = newA[i + newSize / 2][j + newSize / 2];

			B11[i][j] = newB[i][j];
			B12[i][j] = newB[i][j + newSize / 2];
			B21[i][j] = newB[i + newSize / 2][j];
			B22[i][j] = newB[i + newSize / 2][j + newSize / 2];
		}
	}

	vector<vector<int>> M1 = strassenMatrixMultiply(matrixAddition(A11, A22), matrixAddition(B11, B22));
	vector<vector<int>> M2 = strassenMatrixMultiply(matrixAddition(A21, A22), B11);
	vector<vector<int>> M3 = strassenMatrixMultiply(A11, matrixSubtraction(B12, B22));
	vector<vector<int>> M4 = strassenMatrixMultiply(A22, matrixSubtraction(B21, B11));
	vector<vector<int>> M5 = strassenMatrixMultiply(matrixAddition(A11, A12), B22);
	vector<vector<int>> M6 = strassenMatrixMultiply(matrixSubtraction(A21, A11), matrixAddition(B11, B12));
	vector<vector<int>> M7 = strassenMatrixMultiply(matrixSubtraction(A12, A22), matrixAddition(B21, B22));

	vector<vector<int>> C11 = matrixAddition(matrixSubtraction(matrixAddition(M1, M4), M5), M7);
	vector<vector<int>> C12 = matrixAddition(M3, M5);
	vector<vector<int>> C21 = matrixAddition(M2, M4);
	vector<vector<int>> C22 = matrixAddition(matrixAddition(matrixSubtraction(M1, M2), M3), M6);

	vector<vector<int>> result(newSize, vector<int>(newSize));
	for (int i = 0; i < newSize / 2; i++) {
		for (int j = 0; j < newSize / 2; j++) {
			result[i][j] = C11[i][j];
			result[i][j + newSize / 2] = C12[i][j];
			result[i + newSize / 2][j] = C21[i][j];
			result[i + newSize / 2][j + newSize / 2] = C22[i][j];
		}
	}

	// Матрица без лишних нулей
	vector<vector<int>> finalResult(n, vector<int>(n));
	for (int i = 0; i < n; i++) {
		for (int j = 0; j < n; j++) {
			finalResult[i][j] = result[i][j];
		}
	}

	return finalResult;
}

// Функция для умножения матриц классическим способом
vector<vector<int>> classicMatrixMultiply(const vector<vector<int>>& matA, const vector<vector<int>>& matB) {

	if (matA[0].size() != matB.size()) {
		cout << "Невозможно умножить матрицы: неподходящие размеры" << endl;
		return {};
	}

	vector<vector<int>> result(matA.size(), vector<int>(matB[0].size(), 0));

	for (int i = 0; i < matA.size(); i++) {
		for (int j = 0; j < matB[0].size(); j++) {
			for (int k = 0; k < matA[0].size(); k++) {
				result[i][j] += matA[i][k] * matB[k][j];
			}
		}
	}

	return result;
}

// Функция для умножения матриц с использованием алгоритма Винограда
vector<vector<int>> matrixMultiplyVinograd(const vector<vector<int>>& matA, const vector<vector<int>>& matB) {
	int rowsA = matA.size();
	int colsA = matA[0].size();
	int rowsB = matB.size();
	int colsB = matB[0].size();

	if (colsA != rowsB) {
		cout << "Невозможно умножить матрицы: неподходящие размеры" << endl;
		return {};
	}

	vector<vector<int>> result(rowsA, vector<int>(colsB, 0));

	// Предварительные вычисления для оптимизации
	vector<int> rowFactor(rowsA, 0);
	vector<int> colFactor(colsB, 0);
	for (int i = 0; i < rowsA; i++) {
		for (int k = 0; k < colsA - 1; k += 2) {
			rowFactor[i] += matA[i][k] * matA[i][k + 1];
		}
	}
	for (int j = 0; j < colsB; j++) {
		for (int k = 0; k < rowsB - 1; k += 2) {
			colFactor[j] += matB[k][j] * matB[k + 1][j];
		}
	}

	// Основной цикл умножения с использованием предварительных вычислений
	for (int i = 0; i < rowsA; i++) {
		for (int j = 0; j < colsB; j++) {
			result[i][j] = -rowFactor[i] - colFactor[j];
			for (int k = 0; k < colsA - 1; k += 2) {
				result[i][j] += (matA[i][k] + matB[k + 1][j]) * (matA[i][k + 1] + matB[k][j]);
			}
			if (colsA % 2 != 0) {
				result[i][j] += matA[i][colsA - 1] * matB[rowsB - 1][j];
			}
		}
	}

	return result;
}

// Функция для вывода матрицы
void printMatrix(const vector<vector<int>>& mat) {
	for (const auto& row : mat) {
		for (int val : row) {
			cout << val << " ";
		}
		cout << endl;
	}
}

int main(){
	setlocale(0, "");
	cout << "Каким методом Вы хотите перемножить матрицы?" << endl;
	cout << "Классический метод - 1" << endl;
	cout << "Метод Винограда - 2" << endl;
	cout << "Метод Штрассена - 3" << endl;
	int choose = 0;
	double medtime = 0;
	cin >> choose;

	if (choose == 1) {

		vector<vector<int>> matrixA = { {1, 2, 3}, {4, 5, 6}, {7, 8, 9} };
		vector<vector<int>> matrixB = { {9, 8, 7}, {6, 5, 4}, {3, 2, 1} };

		cout << "Матрица A:" << endl;
		printMatrix(matrixA);
		cout << endl;

		cout << "Матрица B:" << endl;
		printMatrix(matrixB);
		cout << endl;

		clock_t start_time = clock();
		vector<vector<int>> result = classicMatrixMultiply(matrixA, matrixB);
		clock_t end_time = clock();
		double cpu_time = (double(end_time - start_time)) / CLOCKS_PER_SEC;
		medtime += cpu_time;
		cout << "Результат умножения классическим методом:" << endl;
		printMatrix(result);
		cout << "Процессорное время: " << cpu_time << endl;
		cout << endl;
	}
	else if (choose == 2) {

		vector<vector<int>> matrixA = { {1, 2, 3}, {4, 5, 6}, {7, 8, 9} };
		vector<vector<int>> matrixB = { {9, 8, 7}, {6, 5, 4}, {3, 2, 1} };

		cout << "Матрица A:" << endl;
		printMatrix(matrixA);
		cout << endl;

		cout << "Матрица B:" << endl;
		printMatrix(matrixB);
		cout << endl;

		clock_t start_time = clock();
		vector<vector<int>> result = matrixMultiplyVinograd(matrixA, matrixB);
		clock_t end_time = clock();
		double cpu_time = (double(end_time - start_time)) / CLOCKS_PER_SEC;
		medtime += cpu_time;
		cout << "Результат умножения с помощью алгоритма Винограда:" << endl;
		printMatrix(result);
		cout << "Процессорное время: " << cpu_time << endl;
		cout << endl;
	}
	else if (choose == 3) {
		vector<vector<int>> matrixA = { {1, 2, 3}, {4, 5, 6}, {7, 8, 9} };
		vector<vector<int>> matrixB = { {9, 8, 7}, {6, 5, 4}, {3, 2, 1} };
		cout << "Матрица A:" << endl;
		printMatrix(matrixA);
		cout << endl;
		cout << "Матрица B:" << endl;
		printMatrix(matrixB);
		cout << endl;

		clock_t start_time = clock();
		vector<vector<int>> result = strassenMatrixMultiply(matrixA, matrixB);
		clock_t end_time = clock();
		double cpu_time = (double(end_time - start_time)) / CLOCKS_PER_SEC;
		medtime += cpu_time;
		cout << "Результат умножения с помощью алгоритма Штрассена:" << endl;
		printMatrix(result);
		cout << "Процессорное время: " << cpu_time << endl;
		cout << endl;
	}
	else {
		cout << "Неверный выбор" << endl;
	}
	return 0;
}
\end{lstlisting}

\newpage
\begin{center}
    \section{Исследовательская часть}
\end{center}

Далее на рисунках представлен результат работы программы.

Результат работы программы умножения матриц $3 \times 3$ классическим способом, приведен на рис.~\ref{img:grap1}. 
\begin{figure}[h!]
    \center{\includegraphics{test1.PNG}}
    \caption{Пример №1}
    \label{img:grap1}
\end{figure}

Результат работы программы умножения матриц $3 \times 3$ методом Винограда, приведен на рис.~\ref{img:grap2}. 
\begin{figure}[h!]
    \center{\includegraphics{test2.PNG}}
    \caption{Пример №2}
    \label{img:grap2}
\end{figure}

\newpage
Результат работы программы умножения матриц $3 \times 3$ методом Штрассена, приведен на рис.~\ref{img:grap3}. 
\begin{figure}[h!]
    \center{\includegraphics{test3.PNG}}
    \caption{Пример №3}
    \label{img:grap3}
\end{figure}

Результат работы программы умножения матриц $6 \times 6$ классическим способом, приведен на рис.~\ref{img:grap4}. 
\begin{figure}[h!]
    \center{\includegraphics{test4.PNG}}
    \caption{Пример №4}
    \label{img:grap4}
\end{figure}

\newpage
Результат работы программы умножения матриц $6 \times 6$ методом Винограда, приведен на рис.~\ref{img:grap5}. 
\begin{figure}[h!]
    \center{\includegraphics{test5.PNG}}
    \caption{Пример №5}
    \label{img:grap5}
\end{figure}

\newpage
Результат работы программы умножения матриц $6 \times 6$ методом Штрассена, приведен на рис.~\ref{img:grap6}. 
\begin{figure}[h!]
    \center{\includegraphics{test6.PNG}}
    \caption{Пример №6}
    \label{img:grap6}
\end{figure}

\newpage
Результат работы программы умножения матриц с различнвми размерами классическим способом, приведен на рис.~\ref{img:grap7}. 
\begin{figure}[h!]
    \center{\includegraphics[scale=0.92]{test7.PNG}}
    \caption{Пример №7}
    \label{img:grap7}
\end{figure}

Классический способ умножения матриц имеет временную сложность $O(n^3)$, где $n$ --- размерность матрицы. Это означает, что время выполнения операции умножения матриц линейно зависит от куба размера матрицы. Как результат, общее количество операций для умножения двух матриц размером $m \times p$ и $p \times n$ составляет $m \cdot p \cdot n$ умножений и $m \cdot n \cdot (p - 1)$ 
сложений.

Умножение матриц с помощью метода Винограда дает больший выигрыш во временной сложности. Она составяет $O(n^{2,3755})$, где $n$ --- размерность матрицы. Результаты подсчета умножений и сложений, где размер матрицы $A$ равен $a \times b$, размер матрицы $B$ равен $b \times c$ и $d = \frac{b}{2}$ наглядно продемонстрированы в таблице~\ref{tab:my_label}.

\begin{table}[h!]
\begin{center}
\caption{Количество операций в методе Винограда}
\label{tab:my_label}
    \centering
    \begin{tabular}{ | m{5cm} | m{5cm}| m{6cm}|} 
  \hline
  Обработка матрицы & Умножение & Сложение \\ 
  \hline
  Предварительная обработка матрицы A & $a \cdot d$ & $a \cdot (d - 1)$\\
  \hline
  Предварительная обработка матрицы B & $c \cdot d$ & $c \cdot (d - 1)$\\
  \hline
  Вычисление матрицы R & $a \cdot c \cdot d$ & $a \cdot c \cdot (2d + d + 1)$\\
  \hline
  Всего & $\displaystyle\frac{abc + ab + bc}{2}$ & $\displaystyle\frac{a(b - 2) + c(b - 2) + ac(3b + 2)}{2}$\\
  \hline
  Сложение двух ненулевых матриц одинакого размера & Ненулевые матрицы M и N & Результирующая ненулевая матрица\\
  \hline
\end{tabular}
\end{center}
\end{table}

Метод Штрассена умножает матрицы за время  $O(n^{\log _{2}7})=O(n^{2.81})$, что даёт выигрыш на больших плотных матрицах. Если обозначить количество операций умножения как $M(n)$, то для обычного метода умножения матриц размера $2^n \times 2^n$ имеем $M(n) = 8^n$. Алгоритм Штрассена снижает количество операций умножения до $M(n) = 7^n$, что означает уменьшение сложности. Он достигает этого за счет уменьшения числа стандартных умножений матриц.

Исходя из вышеперечисленных фатктов, можно сделать вывод, что самым эффективным алогритмом в умножении матриц будет является алгоритм Винограда. Он умножает две матрицы за небольшое количество времени, затрачивает минимальное количетсво памяти и несложен в реализации.

\newpage
\begin{center}
    \section*{Заключение}
    \addcontentsline{toc}{section}{Заключение}
\end{center}
Цель лабораторной работы достигнута. Проведён сравнительный анализ алгоритмов умножения матриц.

В результате выполнения лабораторной работы были выполнены все задачи.
\begin{enumerate}
\item Описаны три алгоритма умножения матриц --- классический, Винограда и Штрассена.
\item Реализованы разработанные алгоритмы.
\item Выполнено тестирование реализации разработанных алгоритмов.
\item Дана оценка сложности алгоритмов.
\item Оценена эффективность каждого алгоритма (по памяти и времени).
\end{enumerate}

\newpage
\begin{center}
\begin{thebibliography}{}
\addcontentsline{toc}{section}{Список используемых источников}
\bibitem{book} Основы современных алгоритмов 2-е дополненное издание / Дж. Макконелл // под ред. С.Кландо, доп. М.В.Ульянова. --- Москва: ТЕХНОСФЕРА, 2004. --- 370 с.
\end{thebibliography}
\end{center}

\end{document}
